\pagebreak
\section{A Wandering Supermassive Black Hole Eating a Star \cite{Saad2026}}} \label{AT2024tvd}
\hrule
\vspace{.5cm}

\begin{list}{-}{}
\item Coveres a paper investigating a \gls{TDE}s dicovered \textit{away} from the center of a galaxy, denoted as AT2024tvd.
\item BH mass \texttt{determines accretion disk behavior}, with more massive BHs having larger, cooler disks
\item This can be measured with the disks' brightness  
\item This was observed with \gls{obs:Pan-STARRS}, \gls{obs:Swift Observatory}, \gls{obs:Pan-STARRS}, and \gls{obs:XMM-Newton}
\item TDEs start with an intial burst in visible and UV which is not very well understood, but after a period of time (few hundred days) the visible and UV light start to come from the accretion disk
\item Used a model called kerrSED, which measures spectral energy distribution (SED), this accounts for disk temp, physical size, spin, viewing angle, and \gls{Comptonization}
\item From this fitting they found that the BH had a mass of \(\sim M_\odot\), classifying it as a \gls{SMBH}  
\item This puts it in a unique spot, as previously detected \gls{TDE}s were located in Ultra-Compact Dwarf Galaxies, small dense collections of stars which orbit larger host galaxies.
\item AT2024tvd has nearly no surrounding stars, suggesting that it is `wandering', having being ejected form a galaxy merger and now slowly drifting towards the galactic center, sheding stars as it goes.   
\end{list}